% answers_short.tex —— 2011 年·辽宁卷·理科 · 简短答案速查
% 由 tools/make_short_answers.py 自 answers.tex 的 \gksolhead/\gkanswer 提取，
% 并与 2011/2011_数学真题_辽宁理_答案要点.md 逐题比对后生成，勿手改（改源 MD 后重跑）。

\begin{gknotice}
\textbf{使用说明}：本节只给最终答案便于快速核对，完整推导见解析版。答案为编写者独立推导，\textbf{非官方标准答案}。
\end{gknotice}

\gkpart{选择题}
\begin{center}\begin{tabular}{ccccccccccccc}
\toprule
题号 & 1 & 2 & 3 & 4 & 5 & 6 & 7 & 8 & 9 & 10 & 11 & 12 \\
\midrule
答案 & B & A & C & D & B & C & A & D & D & B & B & C \\
\bottomrule
\end{tabular}\end{center}

\gkpart{填空题}
\begin{itemize}[leftmargin=2.6em]
  \item[\textbf{13.}] $2$。
  \item[\textbf{14.}] $0.254$。
  \item[\textbf{15.}] $2\sqrt3$。
  \item[\textbf{16.}] $\sqrt3$。
\end{itemize}

\gkpart{解答题}
\begin{itemize}[leftmargin=2.6em]
  \item[\textbf{17.}] (1) $a_n=2-n$；(2) $S_n=\dfrac{n}{2^{n-1}}$
  \item[\textbf{18.}] (1) 证明见解析；(2) $-\dfrac{\sqrt{15}}5$
  \item[\textbf{19.}] (1) $P(X=0)=\dfrac1{70}$，$P(X=1)=\dfrac8{35}$，$P(X=2)=\dfrac{18}{35}$，$P(X=3)=\dfrac8{35}$，$P(X=4)=\dfrac1{70}$；$EX=2$；(2) 甲样本平均数 $400$、方差 $57.25$；乙样本平均数 $412$、方差 $56$；应种植品种乙
  \item[\textbf{20.}] (1) $\dfrac{|BC|}{|AD|}=\dfrac34$；(2) 当 $\dfrac{\sqrt2}2<e<1$ 时存在，此时 $l$ 的横坐标 $t=-\dfrac{a(1-e^2)}{e^2}$；当 $0<e\leqslant\dfrac{\sqrt2}2$ 时不存在
  \item[\textbf{21.}] (1) $a\leqslant0$ 时 $f(x)$ 在 $(0,+\infty)$ 上单调递增；$a>0$ 时在 $\left(0,\dfrac1a\right)$ 上单调递增，在 $\left(\dfrac1a,+\infty\right)$ 上单调递减；(2) 证明见解析；(3) 证明见解析
  \item[\textbf{22.}] (1) A 证明见解析；B $C_1$ 是单位圆，$C_2$ 是椭圆，$a=3$，$b=1$；C 证明见解析；(2) A 证明见解析；B 四边形面积为 $\dfrac25$；C 解集为 $[5-\sqrt3,6]$
\end{itemize}
